\documentclass[12pt, xcolor=dvipsnames,svgnames,x11names]{article}
\usepackage{xcolor}
\usepackage{xspace}
\usepackage{url}
\usepackage[colorlinks=true,bookmarks=false,linkcolor=black,urlcolor=blue,citecolor=black]{hyperref}
\usepackage{fancyhdr}
\usepackage[top=1in,bottom=1.2in,left=1in,right=1in]{geometry}
\usepackage{multicol}
\usepackage{pgfplots}
\usepackage{tikz}
\usetikzlibrary{circuits.logic.US, patterns}
\usepackage{mathpazo}
\usepackage{biolinum} 
\usepackage[all]{tcolorbox}
\usepackage{derivative}
\usepackage{../Lectures/rahul_math}
\renewcommand{\familydefault}{\sfdefault}

% --- Custom Color Palette from Image ---
\definecolor{headerRed}{HTML}{A3403D} % Dark red from image top
\definecolor{patternBg}{HTML}{F2D9D7} % Light pinkish bg from image bottom
\definecolor{binaryText}{HTML}{D1A7A4} % Muted color for binary digits

\usepackage{tikz}
\usetikzlibrary{positioning, arrows.meta, shapes.geometric, calc}

% --- Custom Colors ---
\definecolor{lstmBackground}{RGB}{240, 240, 255}
\definecolor{gateGreen}{RGB}{200, 255, 200}
\definecolor{matrixBlue}{RGB}{135, 206, 250}

% --- Background Pattern Setup ---
\usepackage[pages=all]{background}
\backgroundsetup{
scale=1,
color=black,
opacity=1,
angle=0,
contents={%
  \begin{tikzpicture}[remember picture,overlay]
    % Top Header Bar
    \fill[headerRed] (current page.north west) rectangle ([yshift=-1.2cm]current page.north east);
    % Bottom Binary Background
    \fill[patternBg] (current page.south west) rectangle ([yshift=3.5cm]current page.south east);
    % Decorative Line (from image)
    \draw[headerRed, line width=2pt] ([xshift=1in, yshift=-2.5cm]current page.north west) -- ([xshift=8in, yshift=-2.5cm]current page.north west);
    % Binary Text Overlay (Sample)
    \node[anchor=south, binaryText, font=\ttfamily\scriptsize, inner sep=0pt, yshift=1.5cm] at (current page.south) {
        \begin{tabular}{c@{\hspace{0.5em}}c@{\hspace{0.5em}}c@{\hspace{0.5em}}c@{\hspace{0.5em}}c@{\hspace{0.5em}}c@{\hspace{0.5em}}c}
        0 0 0 0 1 1 1 1 0 0 0 0 1 1 0 1 1 1 0 0 0 1 \\
        1 0 0 1 0 1 0 0 0 1 0 0 1 1 0 0 1 0 1 1 \\
        0 0 1 1 0 0 1 0 1 1 1 0 1 1 0 1 0 1 0 0 0 \\
        The University of Alabama in Huntsville
        \end{tabular}
    };
  \end{tikzpicture}
}
}

\newcommand{\hw}{Classwork 08//Diffusion Process and Gaussian Distribution}
\newcommand{\duedate}{April 07, 2026}

\usetikzlibrary{positioning, arrows.meta, shadows}

% --- Header Styling ---
\makeatletter
\renewcommand{\maketitle}{%
    \vspace*{0.2cm}
    \begin{center}
        \begin{tcolorbox}[
            enhanced,
            colback=white,
            colframe=headerRed,
            arc=0mm,
            title={\textcolor{white}{\bfseries \@title}},
            coltitle=white,
            fonttitle=\bfseries\large,
            attach boxed title to top left={xshift=10pt, yshift=-\tcboxedtitleheight/2},
            boxed title style={sharp corners, colback=headerRed, frame hidden}
        ]
        \large
        \vspace{10pt}
\noindent
\begin{tabular}{@{}p{3.5cm}p{6cm}r@{}}
\textbf{Student Name:} & \underline{\hspace{6cm}} & \textbf{Points:} 30 \\
\end{tabular}

\vspace{10pt}

\noindent
\begin{tabular}{@{}p{3.5cm}p{6cm}r@{}}
\textbf{A Number:} & \underline{\hspace{6cm}} & \textbf{Due Date:} \duedate \\
\end{tabular}
        \end{tcolorbox}
    \end{center}
    \vspace{10pt}
}
\makeatother



\pagestyle{fancy}
\fancyhf{}
\lhead{\small CPE 487/587, Spring 2026}
\rhead{\small \textit{\hw}}
\cfoot{\thepage}
\renewcommand{\headrulewidth}{0pt}

\usepackage{booktabs} % For professional lines


% Define a very light gray for the rows
\definecolor{tablegray}{gray}{0.95}

\begin{document}


\title{\hw}
\maketitle


Consider a random variable
\begin{align*}
\varepsilon \sim \Ncal (0, \Ibf)
\end{align*}

Imagine that $\xbf_0$ is the initial data distribution (the original data distribution). 

Then, in the forward diffusion process, we have 

\begin{align*}
    \xbf_{t+1}  = \sqrt{1 - \beta_t}\xbf_t + \sqrt{\beta_t}\varepsilon
\end{align*}

Here, the recursive relationship collaposes into a single step relative to $\xbf_0$:

\begin{align*}
\xbf_{t+1} = \sqrt{\bar{\alpha}_{t+1}}\xbf_0 + \sqrt{1 - \bar{\alpha}_{t+1}}\varepsilon
\end{align*}

where we have the following relationship:

$\alpha_i = 1 - \beta_i$, and $\bar{\alpha}_{t+1} = \prod_{i=0}^t \alpha_i$


Suppose somebody told you that initial data $\xbf_0$ has distribution $\Ncal(m,s)$ where $s$ is the variance and $m$ is the mean.


\begin{enumerate}
\item Derive the mean $\Ebb[\xbf_{t+1}]$ and the variance $\Vbb(\xbf_{t+1})$ in the form of $\beta_{t+1}$, $m$ and $s$.

\item Also provide the expression when $\beta_{t+1} = 0.1$ for all $t$.

\item Derive Derive $q(\xbf_1 | \xbf_0)$.

\item In Diffusion process, in every step, we make sure that variance of the noisy data at time $t$ doesn't explode. This is called variance preserving method. How does it guide us in choose what $\beta_t$ we can use? How is it related to stochastic differential equation (SDE)?
\end{enumerate}

In doing above derivation, please state what properites of expectations, and Gaussian distributions you are using.

\newpage

\subsection*{Solution}

\begin{enumerate}
\item 

Expectations:

\begin{align*}
\xbf_{t+1} & = \sqrt{\bar{\alpha}_{t+1}}\xbf_0 + \sqrt{1 - \bar{\alpha}_{t+1}}\varepsilon\\
\Ebb[\xbf_{t+1}] & = \Ebb[ \sqrt{\bar{\alpha}_{t+1}}\xbf_0 + \sqrt{1 - \bar{\alpha}_{t+1}}\varepsilon]\\
 & = \sqrt{\bar{\alpha}_{t+1}} \Ebb[\xbf_0] +  \sqrt{1 - \bar{\alpha}_{t+1}} \Ebb[\varepsilon] \quad \text{Property of Expectations}\\
  & = \sqrt{\bar{\alpha}_{t+1}}  m +  \sqrt{1 - \bar{\alpha}_{t+1}}\cdot 0 \quad \textbf{(Given)~~~~} \varepsilon \sim \Ncal (0, \Ibf) \\
& =  \sqrt{\bar{\alpha}_{t+1}}  m
\end{align*}


Variance:

\begin{align*}
\xbf_{t+1} & = \sqrt{\bar{\alpha}_{t+1}}\xbf_0 + \sqrt{1 - \bar{\alpha}_{t+1}}\varepsilon\\
\Vbb[\xbf_{t+1}] & = \Vbb[ \sqrt{\bar{\alpha}_{t+1}}\xbf_0 + \sqrt{1 - \bar{\alpha}_{t+1}}\varepsilon]\\
\Vbb[\xbf_{t+1}] & = ( \sqrt{\bar{\alpha}_{t+1}}) ^2 \Vbb[\xbf_0] + ( \sqrt{1 - \bar{\alpha}_{t+1}} )^2 \Vbb[\varepsilon]\\
& = \bar{\alpha}_{t+1} \Vbb[\xbf_0] + (1 - \bar{\alpha}_{t+1} )\Vbb[\varepsilon]\\
& = \bar{\alpha}_{t+1}  s +  (1 - \bar{\alpha}_{t+1} )\cdot 1  \quad \textbf{(Given)~~~~} \varepsilon \sim \Ncal (0, \Ibf) 
\end{align*}


\item 

\begin{align*}
\beta_{t+1} = 0.1 \forall t\\
\alpha_t = 1 - \beta_t = 0.9\\
\bar{\alpha}_{t+1} = \prod_{i=0}^t \alpha_i = 0.9^{t+1}
\end{align*}


Hence,

\begin{align*}
\Ebb[\xbf_{t+1}] = 0.9 ^{\frac{t+1}{2}}m\\
\Vbb[\xbf_{t+1}] = 0.9^{t+1}s + (1 - 0.9^{t+1} ) 
\end{align*}


\item 

\begin{align*}
q(\xbf_1 | \xbf_0) =?
\end{align*}


Setting $t = 0$, 


\begin{align*}
    \xbf_{1}  = \sqrt{1 - \beta_0}\xbf_0 + \sqrt{\beta_0}\varepsilon\\
\end{align*}

When we condition on a specific value of $\mathbf{x}_0$, the term $\mathbf{x}_0$ is treated as a constant. The only source of randomness in the equation is the Gaussian noise $\varepsilon$. We can find the parameters of $\mathbf{x}_1$ using the properties of linear transformations of Gaussian variables:

\begin{align*}
\mathbb{E}[\mathbf{x}_1 | \mathbf{x}_0] & = \mathbb{E}[\sqrt{1 - \beta_0}\mathbf{x}_0 + \sqrt{\beta_0}\varepsilon]\\
\mathbb{E}[\mathbf{x}_1 | \mathbf{x}_0] & = \sqrt{1 - \beta_0}\mathbf{x}_0 + \sqrt{\beta_0}\mathbb{E}[\varepsilon]
\end{align*}

Since $\mathbb{E}[\varepsilon] = 0$, the mean is:

\begin{align*}
\mathbb{E}[\mathbf{x}_1 | \mathbf{x}_0] & = \sqrt{1 - \beta_0}\mathbf{x}_0 
\end{align*}

For Variance:

$$\mathbb{V}(\mathbf{x}_1 | \mathbf{x}_0) = \mathbb{V}(\sqrt{1 - \beta_0}\mathbf{x}_0 + \sqrt{\beta_0}\varepsilon)$$Because $\mathbf{x}_0$ is constant under the conditioning, its variance is 0.$$\mathbb{V}(\mathbf{x}_1 | \mathbf{x}_0) = (\sqrt{\beta_0})^2 \mathbb{V}(\varepsilon)$$Since $\mathbb{V}(\varepsilon) = \mathbf{I}$, the variance is:$$\mathbb{V}(\mathbf{x}_1 | \mathbf{x}_0)  = \beta_0 \mathbf{I} = \beta_0$$ (if we consider one dimension).

Thus


\begin{align*}
q(\xbf_1 | \xbf_0) \sim \Ncal( \xbf_1;  \sqrt{1 - \beta_0}\mathbf{x}_0 , \beta_0 \mathbf{I})
\end{align*}


\item 

The Variance Preserving (VP) method is a specific design for the forward diffusion process that ensures the latent variables $\mathbf{x}_t$ maintain a constant scale throughout the transition from data to noise. 

In the forward process, the transition is:$$\mathbf{x}_t = \sqrt{1 - \beta_t}\mathbf{x}_{t-1} + \sqrt{\beta_t}\varepsilon$$This is Variance Preserving because if we assume $\mathbf{x}_{t-1}$ has unit variance ($\mathbb{V}(\mathbf{x}_{t-1}) = 1$), then the variance of $\mathbf{x}_t$ remains $1$:$$\mathbb{V}(\mathbf{x}_t) = (\sqrt{1 - \beta_t})^2 \cdot 1 + (\sqrt{\beta_t})^2 \cdot 1 = 1 - \beta_t + \beta_t = 1$$

Contrast this with a simple random walk (like Variance Exploding methods), where you simply add noise: $\mathbf{x}_t = \mathbf{x}_{t-1} + \sigma_t\epsilon$. In that case, the variance would grow to infinity as $t$ increases.

But the VP formulation acts as a guardrail for the generative model.

Because the total variance is kept near $1$, the model doesn't have to deal with exploding numerical values. The data is effectively squashed toward the origin as noise is added, ensuring $\mathbf{x}_T$ always lands within the neat confines of a standard normal distribution $\mathcal{N}(0, \mathbf{I})$.

VP diffusion can be viewed as a discretization of the following SDE:$$d\mathbf{x} = -\frac{1}{2}\beta(t)\mathbf{x} \, dt + \sqrt{\beta(t)} \, d\mathbf{w}$$

Here, the first term ($-\frac{1}{2}\beta(t)\mathbf{x}$) is a drift term that constantly pulls the distribution back toward the origin, preventing it from wandering off, while the second term adds the diffusion.

\end{enumerate}



\end{document}